% A local honeycomb neighborhood with the usual negatively sloping
% reservoir transitions and positively sloping interdot transitions.
\begingroup
\small
\setlength{\unitlength}{1mm}
\begin{picture}(150,87)
 \put(15,8){\vector(1,0){122}}
 \put(15,8){\vector(0,1){73}}
 \put(137,4){\makebox(10,7){$V_1$}}
 \put(5,78){\makebox(10,7){$V_2$}}
 \put(66,35){\makebox(25,8){$(1,1)$}}
 \put(32,26){\makebox(25,8){$(0,1)$}}
 \put(53,11){\makebox(25,8){$(1,0)$}}
 \put(90,18){\makebox(25,8){$(2,0)$}}
 \put(106,44){\makebox(25,8){$(2,1)$}}
 \put(83,62){\makebox(25,8){$(1,2)$}}
 \put(28,62){\makebox(25,8){$(0,2)$}}
 \linethickness{0.9pt}
 \color{blue!65!black}
 % Central cell A(51,29), B(83,21), C(103,37),
 % D(95,61), E(63,69), F(43,53).
 \qbezier(51,29)(67,25)(83,21)
 \qbezier(103,37)(99,49)(95,61)
 \qbezier(95,61)(79,65)(63,69)
 \qbezier(43,53)(47,41)(51,29)
 % Reservoir-transition edges extending to adjacent cells.
 \qbezier(83,21)(85,15)(87,9)
 \qbezier(103,37)(118,33.25)(133,29.5)
 \qbezier(63,69)(61,75)(59,81)
 \qbezier(43,53)(31.5,55.875)(20,58.75)
 \color{orange!85!black}
 \qbezier(83,21)(93,29)(103,37)
 \qbezier(63,69)(53,61)(43,53)
 \qbezier(51,29)(40,20.2)(29,11.4)
 \qbezier(95,61)(107,70.6)(119,80.2)
 \color{black}
 \put(51,29){\circle*{2}}
 \put(83,21){\circle*{2}}
 \put(103,37){\circle*{2}}
 \put(95,61){\circle*{2}}
 \put(63,69){\circle*{2}}
 \put(43,53){\circle*{2}}
\end{picture}
\endgroup
